\section{Resultados Esperados}
\label{sec:resultados_esperados}

\begin{itemize}

  \item \textbf{Modelo calibrado.} Um modelo SEIR-SEI ajustado à primeira onda de chikungunya no Recife,
  com estimativas e intervalos de confiança para o número de reprodução, a fração de notificação e a
  sazonalidade do vetor.

  \item \textbf{Imunidade remanescente.} Estimativas da fração da população ainda suscetível na cidade e,
  com o modelo por grupos, em cada distrito sanitário ou estrato.

  \item \textbf{Cenários de controle.} Comparação entre redução de criadouros e aumento da mortalidade dos
  mosquitos adultos, com diferentes eficácias e momentos de início, expressa em casos evitados.

  \item \textbf{Implementação computacional aberta.} Repositório público com código, dados agregados,
  notebooks e o script que regenera todos os resultados conferindo-os contra a teoria.

  \item \textbf{Painel de monitoramento.} Um painel que mostra a série observada, o ajuste do modelo e os
  cenários, atualizado a cada novo arquivo anual do portal de dados abertos.

  \item \textbf{Impacto institucional.} Um relatório técnico que possa ser apresentado à vigilância
  epidemiológica do Recife, e material didático para a disciplina de Epidemiologia Computacional, que
  mostra o caminho do SIR ao modelo com vetor usando dados da própria cidade.

\end{itemize}
